Cross-regional delay surveys converge on a recurring causal taxonomy even when ranked importance shifts with institutional setting. Assaf and Al-Hejji reported change-order dominance and widespread lateness in large Saudi projects~\cite{Assaf2006}. Sambasivan and Soon linked Malaysian delay causes to secondary effects including disputes and abandonment~\cite{Sambasivan2007}. Doloi \textit{et al.} emphasized managerial commitment and site coordination~\cite{Doloi2012}. Nasir \textit{et al.}'s ERIC-S framework and Kim and Reinschmidt's Bayesian duration updating illustrate structured taxonomies and sequential belief revision~\cite{Nasir2003,Kim2016}. Flyvbjerg's optimism-bias analyses and Love \textit{et al.}'s transport-overrun synthesis argue that forecast error is systematic rather than purely aleatory~\cite{Flyvbjerg2009,Love2015}. Those surveys motivate budgeted inspection without specifying a numeric policy objective.

Bayesian network applications then made causal structure explicit. Luu \textit{et al.} quantified schedule risk with BBNs~\cite{Luu2009}. Ji and Leiringer developed a UK delay-avoidance predictive BBN~\cite{Ji2021}; Comu \textit{et al.} quantified marine-project delay risk~\cite{Comu2023}; Khodabakhshian \textit{et al.} combined elicitation with structural learning and reported post-augmentation accuracy near $0.85$~\cite{Khodabakhshian2025}; Zhang \textit{et al.} coupled structure learning with Monte Carlo schedule inference~\cite{Zhang2021BNMCS}. Related construction-safety BBNs demonstrate evidence updating under field observations~\cite{Wu2015,Wang2016,Leu2012}. Across these papers, expert structure plus CPT quantification can yield usable risk scores, yet sample sizes, hold-out designs, and decision metrics vary, and sequential inspection under a hard budget is seldom the evaluation target.

Outside construction, VoI is a mature inspection-planning instrument. Howard's formulation defines expected utility gain from an observation before it arrives~\cite{Howard1966}. Memarzadeh and Pozzi analyzed sequential VoI for component inspection~\cite{Memarzadeh2016}; Malings and Pozzi treated spatially distributed sensing~\cite{Malings2016}; Song \textit{et al.} specialized VoI metrics to binary systems~\cite{Song2021}; Yuan \textit{et al.} valued inspection programs under Bayesian degradation models~\cite{Yuan2021}; Fauriat and Zio optimized aperiodic inspection with VoI-driven maintenance~\cite{Fauriat2020}. These studies establish both the mathematics and the habit of reporting when VoI differentials vanish---a habit construction delay papers have rarely adopted.

Probabilistic evaluation supplies a second tradition that construction delay ML still underuses. Brier's score and Murphy's decomposition connect probabilistic forecast quality to calibration and refinement~\cite{Brier1950,Murphy1973}. Gneiting and colleagues clarify sharpness, calibration, and strictly proper scoring~\cite{Gneiting2007,GneitingRaftery2007}. Saito and Rehmsmeier argue for precision--recall analysis under imbalance~\cite{Saito2015}; Hand discusses ROC interpretation pitfalls~\cite{Hand2009}; Vickers and Elkin introduce decision-curve net benefit~\cite{Vickers2006}. Cohen's $\kappa$ and Matthews' correlation provide chance-corrected agreement summaries~\cite{Cohen1960,Matthews1975}. Construction AI reviews document rapid method adoption~\cite{Darko2020,PanYZhang2021,Abioye2021,Regona2024} and classifier bake-offs report strong accuracies~\cite{Egwim2021,Tagharobi2025}, yet calibrated probability reporting remains uneven.

Hybrid and fuzzy risk models occupy a parallel branch. Ebrat and Ghodsi used adaptive neuro-fuzzy inference for construction project risk~\cite{Ebrat2014}; Islam \textit{et al.} reviewed fuzzy hybrid techniques~\cite{Islam2017}; Qazi \textit{et al.} examined project complexity and risk paths~\cite{Qazi2016}; Zou \textit{et al.} reviewed BIM-linked risk management~\cite{Zou2017}. Woldesenbet \textit{et al.} stressed provenance and integration quality for highway infrastructure data~\cite{Woldesenbet2016}. These strands reinforce preferring original open-data endpoints and measuring free-text reason-mapping error rather than assuming it away. Collectively, the literature implies that methodological novelty alone is insufficient: leakage-safe splits, proper scoring, and policy-identifiability under realistic sparsity determine whether a tool can support inspection decisions.
